

\section{Generic Parser Algorithm with Effects and Evidence}
\label{sec:prefix-parsing-algorithm}

The core algorithm is a language-parametric prefix parser. It can serve as a static analysis tool, a completion engine, or a full type checker.

\subsection{Language-Parametric Interface}

The framework is parametric over SPGs: a specification defines productions, binding annotations, and semantic rules. The engine checks well-formedness, builds binding maps as in Section~\ref{binding}, and then exposes the same prefix parser for every language instance.

\subsection*{Obligations}


Semantic rules are not resolved only over finished trees. Instead, they are computed at parse time and used for branch pruning through an \emph{obligation} system, which is the parser-time realization of the binding-resolution map $\beta$ of Definition~\ref{def:beta-resolution}. The behavior is structurally equivalent, but divides the evaluation step into two operations : (i) obligation instantiation and (ii) final resolution. Pruning is done thanks to the binding information gathered in (i).

\begin{definition}[Obligation and store]
\label{def:obligation-store}
An \emph{obligation store} $\Omega$ records the evidence expected from bound children. Each obligation has the form
\[
  \omega_b ~=~ (b,~v_b,~E_b),
\]
where $b$ is the binding name, $v_b$ is terminal or span evidence, and $E_b$ is domain evidence. An obligation $\omega_b$ with both $v_b = \bot$ and $E_b = \bot$ is the parser-side encoding of $\beta(\nu,b) = \mathsf{Pending}(i_b)$; resolving a terminal child stores terminal evidence and yields $\beta(\nu,b) = \mathsf{Resolved}(\mathsf{Term}(\cdot))$; resolving a non-terminal child stores span and exported domain evidence and yields $\beta(\nu,b) = \mathsf{Resolved}(\nu_{c_{i_b}})$.
\end{definition}

The purpose of obligations is to provide an interface for rule evaluation that behaves like recursive descent, while integrating with the parser to reduce complexity and enable pruning. Equivalently, the store materializes the $\mathsf{Pending}$-to-$\mathsf{Resolved}$ transitions of $\beta$ along the dot.

%
%\begin{lemma}[Obligation locality]
%\label{lem:obligation-locality}
%For a binding-regular grammar, the obligations projected to position~$i$ of
%alternative $p^n_a$ are exactly
%\[
%  \mathsf{step}(\Omega, i, a)
%  \;=\;
%  \{\, \omega_b \in \Omega \mid \mathsf{pos}_{n,a}(b) = i \,\}.
%\]
%Resolving position~$i$ may modify only these obligations.
%\end{lemma}
%
%\begin{proof}
%By binding uniqueness (Property~\ref{prop:binding-uniqueness}), each binding~$b$
%has a unique position $\mathsf{pos}_{n,a}(b)$ in $p^n_a$. Position~$i$ can
%supply evidence only for bindings whose unique position is~$i$.
%\end{proof}

\subsection{Parsing}
\label{sec:engine-parsing}

A \emph{prefix parser} for an SPG $G$ is a function
\[
\Psi_G : \Sigma^* \to \mathcal F_\partial \cup \{\texttt{reject}\}.
\]

\noindent It is defined by a runtime state $\gamma$ that includes:
\begin{itemize}
  \item The input segmentation $x=\mathsf{seg}_{\mathcal L}(s)=(x_0,\ldots,x_{m-1})$
  \item The agenda $\mathcal A$ of items to process
  \item The chart $C$ mapping recognized nonterminal spans to node identifiers
  \item The waiter relation $W$ mapping nonterminal spans to paused parent items
  \item The set $R$ of known end positions for each nonterminal and segment index
  \item The end-of-input frontier $F$ of items awaiting EOF closure
  \item The packed parse arena $H$
\end{itemize}

\noindent A live item is
\[
\eta=\langle n,p,k,[i,j),\Gamma,\Omega,\chi \rangle
\]
where $n$ is the non-terminal being recognized, $p=p^n_a$ is one of its productions, $k$ is the dot position, $[i,j)$ is the consumed segment span, $\Gamma$ is the semantic context at the dot, $\Omega$ is the current obligation store, and $\chi$ is the ordered list of accepted child references and terminal evidence.

\[
\begin{array}{c}
p =
\underbrace{\alpha_1[b_1]\ \cdots\ \alpha_k[b_k]}_{\text{recognized}}
\ \bullet\
\underbrace{\alpha_{k+1}[b_{k+1}]\ \cdots\ \alpha_n[b_n]}_{\text{remaining}}
\end{array}
\]


Each arena node is the evidence summary $\nu$ of Definition~\ref{def:evidence-summary}, viewed as a packed handle for a family of parses and extended with resolved domain evidence $\tau$ that makes already evaluated semantic resolution available and a right-bound effect $\nabla$ that exposes how the node will modify the semantic context of right siblings.
\[
\nu=\langle n,[i,j),\rho,\tau,\nabla,\beta \rangle
\]
with $\rho \in \{\mathsf{Exact},\mathsf{Prefix},\mathsf{Extensible}\}$, $\tau$ the exported semantic evidence of the active constraint domain, $\nabla$ an optional right-bound effect, and $\beta$ the parent-visible binding-resolution map (Definition~\ref{def:beta-resolution}), where $\beta(\nu,b)$ is realized by the obligation $\omega_b \in \Omega$ of Definition~\ref{def:obligation-store}. Exact nodes are closed on the current input and may export $\nabla$. Prefix nodes are syntactically unfinished at the right frontier. Extensible nodes are fully recognized on the current input but contain right-frontier terminal evidence that may still grow. Prefix and Extensible nodes never export effects.



\subsubsection{Semantic Interface}

The parser uses the following operations. Mathematically, mutating operations are written as functional updates. They are partial: failure means that the branch is discarded.
\[
\begin{array}{rcl}
\mathsf{create}(G,p,r)
&:&
\text{creates production-local obligations for production }p\text{ at root }r,
\\[0.25em]

\mathsf{step}(\Omega,k,a)
&:&
\text{projects obligations through child position }k\text{ of alternative }a,
\\[0.25em]

\mathsf{prune}(\Omega,n,G)
&:&
\text{returns the alternatives of }n\text{ compatible with }\Omega,
\\[0.25em]

\mathsf{for}(\Omega,a)
&:&
\text{keeps only obligations compatible with seeded alternative }a,
\\[0.25em]

\mathsf{at}(\Omega,k)
&:&
\text{re-roots obligations at child position }k,
\\[0.25em]

\mathsf{descend}(p,b,\Gamma,\Omega)
&:&
\text{computes the context used to enter child binding }b,
\\[0.25em]

\mathsf{resolve}_T(\Omega,k,a,\mathsf{Term}(x,\zeta,r))
&:&
\text{stores terminal evidence for obligations at }(k,a),
\\[0.25em]

\mathsf{resolve}_N(\Omega,k,a,\nu)
&:&
\text{stores span and domain evidence from accepted child node }\nu,
\\[0.25em]

\mathsf{finalize}(p,\Gamma,\Omega,\rho)
&:&
\text{evaluates the production rule when finishing a node},
\\[0.25em]

\mathsf{apply}(\Gamma,\nabla)
&:&
\text{applies an exported right-bound effect}.
\end{array}
\]
We assume all semantic operations are denotational filters: defined outputs denote exactly the satisfying refinements, and semantic rejection happens only on empty denotation.

\subsubsection{Push Transitions}

All parser table updates are written as
\[
\gamma \downarrow_{\alpha}(y),
\]
where $\alpha$ names the table or helper being updated and $y$ is the payload supplied to that update. The update labels are
\[
\alpha ::=
\mathsf{process}
\mid
\mathsf{waiters}
\mid
\mathsf{frontier}
\mid
\mathsf{node}
\mid
\mathsf{record}
\mid
\mathsf{complete}.
\]
The intended payloads are:
\[
\begin{array}{rcl}
\gamma \downarrow_{\mathsf{process}}(\eta)
&:&
\text{interns a process task in the agenda, subject to the item key},
\\[0.25em]

\gamma \downarrow_{\mathsf{waiters}}(n,j,\eta)
&:&
\text{interns the paused parent item }\eta\text{ in }W(n,j),
\\[0.25em]

\gamma \downarrow_{\mathsf{frontier}}(\eta)
&:&
\text{appends }\eta\text{ to }F,
\\[0.25em]

\gamma \downarrow_{\mathsf{node}}(\nu,\pi)
&:&
\text{interns a node and packed alternative in }H,
\\[0.25em]

\gamma \downarrow_{\mathsf{record}}(c)
&:&
\text{records a completed node in }C\text{ and its end position in }R,
\\[0.25em]

\gamma \downarrow_{\mathsf{complete}}(c)
&:&
\text{appends a completion task and wakes matching waiters.}
\end{array}
\]
When a push allocates a fresh arena handle, we write the handle above the push:
\[
\gamma {\downarrow}^h_{\mathsf{node}}(\nu,\pi).
\]

\subsubsection{Parse Inference}

The following rules summarize the semantic fields added to the standard prefix-Earley transition system. They use the push transition explicitly; standard side conditions such as deduplication and completed-dot checks are the side conditions of the corresponding push action.

We write $\Omega_1\oplus\Omega_2$ for extending one obligation store with another.

\[
\frac{
  q=p^n_a \in Q \subseteq P
  \quad
  \Omega^{\mathsf{in}}_q=\mathsf{for}(\Omega,a)
  \quad
  \Omega_q=
  \mathsf{create}(G,q,\mathsf{root}(\Omega^{\mathsf{in}}_q))
  \oplus
  \Omega^{\mathsf{in}}_q
}{
\gamma \downarrow_{\mathsf{process}}
  (\langle n,q,0,[j,j),\Gamma,\Omega_q,\epsilon\rangle)
}
\quad\textsc{Seed}
\]
Seed is parameterized by a set of candidate productions \(Q\), an input
position \(j\), an incoming context \(\Gamma\), and an incoming obligation store
\(\Omega\). The initial call is $Q=A(S)$, $j=0$, $\Gamma=\Gamma_0$, and $\Omega=\mathsf{empty}$.

\[
\frac{
  \begin{array}{cr}
  &\eta=\langle n,p^n_a,k,[i,j),\Gamma,\Omega,\chi\rangle
  \quad
  p^n_a[k]=n'[b]
  \quad
  n'\in N
  \quad \\[0.30em]
  &\Gamma_c=\mathsf{descend}(p^n_a,b,\Gamma,\Omega)
  \quad
  \Omega^\dagger=\mathsf{step}(\Omega,k,a)
  \quad
  \Omega_c=\mathsf{at}(\Omega^\dagger,k)
  \end{array}
}{
\gamma \downarrow_{\mathsf{waiters}}(n',j,\eta)
\qquad
\forall q \in \mathsf{prune}(\Omega^\dagger,n',G),\quad
\gamma \downarrow_{\mathsf{process}}(\eta_q)
}
\quad\textsc{ProcessNonterminal}
\]
where each $\eta_q$ is the child item produced by \textsc{Seed} with $P=\{q\}$, position $j$, context $\Gamma_c$, and parent obligations $\Omega_c$. Thus pruning is applied before any child production is seeded, and $W(n',j)$ stores the paused parent item rather than the predicted child.

\[
\frac{
\begin{array}{c}

  \eta=\langle n,p^n_a,k,[i,j),\Gamma,\Omega,\chi\rangle
  \quad
  p^n_a[k]=t[b]
  \quad
  t\in T 
  \quad
  \mathsf{match}(t,x_j)=(\zeta,r) \\[0.25em]
  \quad
  \zeta\in\{\mathsf{Exact},\mathsf{Extensible}\} 
  \quad
  T_j=\mathsf{Term}(x_j,\zeta,r)
  \quad
  \Omega'=\mathsf{resolve}_T(\Omega,k,a,T_j)
\end{array}
}{
  \gamma \downarrow_{\mathsf{process}}
  (\langle n,p^n_a,k+1,[i,j+1),\Gamma,\Omega',\chi\cdot T_j\rangle)
}
\quad\textsc{ConsumeReady}
\]

\[
\frac{
\begin{array}{c}
  \eta=\langle n,p^n_a,k,[i,j),\Gamma,\Omega,\chi\rangle
  \quad
  p^n_a[k]=t[b]
  \quad
  t\in T
  \quad
  j=m-1 \\[0.25em]
  \quad
  \mathsf{match}(t,x_j)=(\mathsf{Prefix},r)
  \quad
  T_j=\mathsf{Term}(x_j,\mathsf{Prefix},r)
  \quad
  \Omega'=\mathsf{resolve}_T(\Omega,k,a,T_j)
\end{array}
}{
  \gamma \downarrow_{\mathsf{frontier}}
  (\langle n,p^n_a,k+1,[i,m),\Gamma,\Omega',\chi\cdot T_j\rangle)
}
\quad\textsc{ConsumePrefix}
\]

\[
\frac{
  \eta=\langle n,p,k,[i,m),\Gamma,\Omega,\chi\rangle
  \quad
  p[k]\in T
}{
  \gamma \downarrow_{\mathsf{frontier}}(\eta)
}
\quad\textsc{ConsumeEof}
\]

Let $\rho(\eta)$ be the status selected by $\mathsf{finish}$:
\[
\rho(\eta)=
\begin{cases}
\mathsf{Exact} & \text{all children are complete and none is open},\\
\mathsf{Extensible} & \text{all children are complete and at least one is open},\\
\mathsf{Prefix} & \text{otherwise}.
\end{cases}
\]
Prefix and Extensible statuses force the exported effect to $\bot$.

\[
\frac{
  \begin{array}{cc}
  \eta=\langle n,p,k,[i,j),\Gamma,\Omega,\chi\rangle
  \quad
  \mathsf{finishable}(\eta)
  \quad
  \rho=\rho(\eta)
  \quad
  \mathsf{finalize}(p,\Gamma,\Omega,\rho)=(\tau,\nabla,\spadesuit) \\[0.25em]
  \quad
  \nabla_\rho=
  \begin{cases}
    \nabla & \rho=\mathsf{Exact},\\
    \bot & \rho\in\{\mathsf{Prefix},\mathsf{Extensible}\},
  \end{cases} 
  \quad
  \nu=\langle n,[i,j),\rho,\tau,\nabla_\rho,\mathsf{bindings}(\Omega)\rangle
  \quad
  \pi=(p,\chi)
  \end{array}
}{
  \gamma {\downarrow}^h_{\mathsf{node}}(\nu,\pi)
  \qquad
  \gamma \downarrow_{\mathsf{complete}}(\langle n,i,j,h\rangle)
}
\quad\textsc{Finish}
\]
Here $\mathsf{finishable}(\eta)$ means either that the main $\mathsf{process}$ loop has reached the end of the production, or that $\eta$ has been identified as partial from the EOF frontier. The value $\spadesuit$ is the semantic-completeness flag status, corresponding to the $\mathit{eval}$ \ref{def:constraint-domain} return values. It is used to trigger branch abandon if it is $\mathsf{Lost}$. Finalize acts as a sort of "closing" for a  $\mathit{eval}$ call.

\[
\frac{
\begin{array}{cr}
  c=\langle n',j,\ell,h\rangle
  \quad
  \eta=\langle n,p^n_a,k,[i,j),\Gamma,\Omega,\chi\rangle\in W(n',j)
  \quad
  H[h]=\nu=\langle n',[j,\ell),\rho,\tau,\nabla,\beta\rangle \\[0.25em]
  \quad
  \Omega'=\mathsf{resolve}_N(\Omega,k,a,\nu)  
  \quad
  \Gamma'=\begin{cases}
    \mathsf{apply}(\Gamma,\nabla) & \nabla\neq\bot, \\
    \Gamma & \nabla=\bot,
  \end{cases}
  \quad
  \eta'=\mathsf{resume}(\eta,h)
\end{array}
}{
  \gamma \downarrow_{\mathsf{process}}(\eta')
}
\quad\textsc{CompleteResume}
\]
The completion $c$ is pushed once, by \textsc{Finish}, with
\[
\gamma \downarrow_{\mathsf{complete}}(c).
\]
Then the resume step above is repeated for each waiter in $W(n',j)$ for which $\mathsf{resume}$ is defined.
\medskip

After the main agenda is saturated, we do \emph{frontier lifting} on $F$ by repeatedly applying \textsc{Finish} and \textsc{CompleteResume} to frontier items and their resumed parents. It does not call $\mathsf{process}$, $\mathsf{seed}$, or $\mathsf{consume}$, but rather it only finishes frontier summaries and propagates them through the existing waiter network. This is a central element to what makes our parser able to process \emph{prefixes}

The parser succeeds when saturation and EOF closure produce a root node
\[
\nu=\langle S,[0,m),\rho,\tau,\nabla,\beta\rangle
\quad
\rho\in\{\mathsf{Exact},\mathsf{Prefix},\mathsf{Extensible}\}.
\]
If no such root node is produced, the parser returns \texttt{reject}.

\subsection{Context and Effect Flow}
\label{sec:context-flow}

The context stored in an item is the context at the dot. It is obtained by combining two forms of flow: top-down context passed into a child, and right-bound effects exported by exact children to their right siblings.

\paragraph{Downward context}

When the parser predicts a child $p[k]=n'[b]$, it computes the child context by
\[
\Gamma_{\mathsf{child}}
=
\mathsf{descend}(p,b,\Gamma,\Omega).
\]
This is a downward flow: the child is parsed under the context selected by the parent rule. For example, the body of a lambda is predicted under $\Gamma[x:\tau]$, but this extension is local to that child. This property is inherited from recursive descent, where it is more apparent.

\paragraph{Right-bound flow}

Right-bound flow happens after a child is completed. If exact children to the left of the dot export effects
\[
\nabla_1,\ldots,\nabla_q,
\]
then the context at the dot is
\[
\Gamma_q=(\nabla_q\circ\cdots\circ\nabla_1)(\Gamma_0).
\]
Accepting the next exact child with export $\nabla_{q+1}$ updates the parent context to
\[
\Gamma_{q+1}
=
\nabla_{q+1}(\Gamma_q)
=
(\nabla_{q+1}\circ\nabla_q\circ\cdots\circ\nabla_1)(\Gamma_0).
\]
Non-exact nodes do not participate in this composition: Prefix and Extensible nodes export $\bot$, so resuming such a child leaves the parent context unchanged. Thus top-down flow determines the context used inside a child, while right-bound flow determines the context seen by later siblings. In the typing-domain instance, the effects $\nabla_i$ are context transforms; other domains may instantiate them differently.

\subsection{Soundness}
\label{sec:parser-soundness}

The parser is used as a prefix decision procedure. The property we need is that every accepted prefix can still be completed into a closed semantic parse. 

The invariant is simple: accepted means not dead. Safe pruning lets the parser work over an infinite continuation tree because it only cuts branches that cannot hold in the semantic domain. We use two invariants of the rules above. First, after erasing semantic fields, the rules form the standard Earley prefix parser for the syntactic projection of $G$. Second, semantic operations are exact local filters: a semantic operation is defined exactly when the evidence graph carried by the current item or node is not lost. In particular, exact finalization requires a closed satisfying graph, while non-exact finalization requires a live graph in the sense of Definition~\ref{def:constraint-domain}.

\begin{theorem}[Completability soundness]
\label{thm:completability-soundness}
Let $G\in \spg(D)$ be productive by Assumption~\ref{ass:spg-productivity}. If
\[
  \Psi_G(s) \neq \mathsf{reject},
\]
then there exists $s' \in \Sigma^\ast$ such that $\Psi_G(ss')$ produces an $\mathsf{Exact}$ root node.
\end{theorem}

\begin{proof}
Since $\Psi_G(s) \neq \mathsf{reject}$, saturation on $s$ produced a root node
\[
  \nu = \langle S,[0,m),\rho,\tau,\nabla,\beta\rangle
\]
with $\rho \in \{\mathsf{Exact},\mathsf{Prefix},\mathsf{Extensible}\}$.

If $\rho = \mathsf{Exact}$, the root is already closed on the current input, so taking $s'=\varepsilon$ proves the result.

Suppose instead that $\rho \in \{\mathsf{Prefix},\mathsf{Extensible}\}$. This root was produced by non-exact finalization. By the semantic side condition of $\textsc{Finish}$, the $\spadesuit$ flag of $\mathit{eval}_{\mathrm{impl}}$ on $\nu$ is non-$\mathsf{Lost}$, and by Theorem~\ref{thm:typing-realizable} this matches $\mathit{eval}(\mathcal G(s))$ on the same evidence summary. Since $\rho \neq \mathsf{Exact}$, at least one frontier vertex remains in the sense of Definition~\ref{def:open-closed}, so $\mathit{eval}(\mathcal G(s)) = \mathsf{Live}$. By Definition~\ref{def:constraint-domain}, this means
\[
\exists s'\in\Sigma^\ast.\;\mathcal G(s\cdot s') \in \mathsf{Closed},
\]
and by Definition~\ref{def:denotation}, such an $s'$ is precisely a witness $s' \in \mathcal D(\mathcal G(s))$.

It remains only to justify that the parser will actually produce the closed root on $ss'$. The suffix $s'$ closes the branch summarized by $\nu$. Syntactically, the erased inference system is the Earley prefix parser, so the corresponding syntactic continuation is admitted. Semantically, each operation on that branch is an exact filter, so following a satisfying continuation cannot make a required semantic operation fail. Since the system is saturated, all items and nodes on this closed branch are eventually inserted. In particular, $\Psi_G(ss')$ produces an $\mathsf{Exact}$ root node.
\end{proof}

\begin{lemma}[Prefix completeness of saturation]
\label{lem:prefix-completeness}
Let $G \in \mathsf{SPG}(D)$. If $\Psi_G(sr)$ produces an $\mathsf{Exact}$ root node, then
\[
  \Psi_G(s) \neq \mathsf{reject}.
\]
\end{lemma}

\begin{proof}
Take a saturated derivation of the exact root over $sr$, and cut it at the input frontier after $s$. By prefix-awareness of the segmentation policy, closed interior segments are unchanged by this cut; only the final segment may become open.

The syntactic projection of the cut derivation is therefore a valid prefix Earley derivation for $s$. Semantic evidence already to the left of the cut is unchanged. Evidence crossing the cut becomes prefix evidence, and it is retained because the complete derivation over $sr$ gives an actual suffix closing it.

Moreover, semantic graphs grow monotonically with input:
\[
  \mathcal E(s) \subseteq \mathcal E(sr)
  \qquad\text{and}\qquad
  \mathcal C(s) \subseteq \mathcal C(sr).
\]
This holds because each parser transition only adds vertices and edges: $\textsc{Seed}$ and $\textsc{ProcessNonterminal}$ insert fresh items, $\textsc{Finish}$ interns a node $\nu$, and $\textsc{CompleteResume}$ flips $\beta(\nu,b)$ from $\mathsf{Pending}(i_b)$ to $\mathsf{Resolved}(\nu_{c_{i_b}})$ (Definition~\ref{def:beta-resolution}) without removing the binding-reference vertex. The exact root over $sr$ has a closed satisfying graph. Restricting that closed graph to the evidence and constraints already present at $s$ gives a satisfying grounding for the prefix graph. Hence none of the semantic operations in the cut derivation is lost. 
Since the inference system is saturated, the corresponding exact or non-exact root node is inserted when parsing $s$. Therefore $\Psi_G(s) \neq \mathsf{reject}$.
\end{proof}

\begin{theorem}[Prefix monotonicity]
\label{thm:prefix-monotonicity}
Let $G \in \mathsf{SPG}(D)$ be productive by \ref{ass:spg-productivity}. If
\[
  \Psi_G(s) \neq \mathsf{reject},
\]
then for every prefix $s''$ of $s$,
\[
  \Psi_G(s'') \neq \mathsf{reject}.
\]
\end{theorem}

\begin{proof}
Let $s=s''r$. By Theorem~\ref{thm:completability-soundness}, there exists $t \in \Sigma^\ast$ such that
\[
  \Psi_G(st)=\Psi_G(s''rt)
\]
produces an $\mathsf{Exact}$ root node. Applying Lemma~\ref{lem:prefix-completeness} to the complete input $s''rt$ gives
\[
  \Psi_G(s'') \neq \mathsf{reject}.
\]
\end{proof}
